\documentclass{article}
\usepackage{amsthm}

\usepackage[utf8]{inputenc}     % pdfLaTeX
\usepackage[T1]{fontenc}
\usepackage[magyar]{babel}

\title{Ázsia}
\author{Haydn Lorenzo}
\date{Legutóbb frissítve: 2026. Szeptember 17.}

\begin{document}

\maketitle
\tableofcontents
\newpage

\section{Ázsia bomborzata}
\subsection{Ősföldek}
\begin{itemize}
    \item Augara: Közép-Szibériai-fennsík: fedett ősmaszívum
    \item Kínai:\begin{itemize}
        \item É: üledékkel fedett táblás vidék
        \item D: Dél-Kínai-hegyvidék
    \end{itemize}
    \item Dekkán-ősföld: \begin{itemize}
        \item Ny: bazalt takaró
        \item: K: fedetlen
    \end{itemize}
    \item Arab-ősföld: \begin{itemize}
        \item tengeri üledékkel fedett: kőolaj, földgáz
    \end{itemize}
\subsection{Óidei röghegység}
\begin{itemize}
    \item Urál: rögökre darabolódott, tömbökből álló hegyvidék
    \item Tien-San: Óidőben kialakult röghegység lekopott, Himalája keletkezésekor újra felgyűrődött
    \item magashegységi formakincs
\end{itemize}
\end{itemize}

\end{document}